% ------------------------------------------------------------
% A6 Firewall — Dual-Algebra Boundary Conditions (Module 1)
% Author: Borealis S. Hedling
% Altitude: A6.5 (Boundary Layer)
% Purpose: Define the structural firewall separating Euclidean
%          dual-algebra from non-Euclidean non-dual algebra.
% ------------------------------------------------------------

\section{The A6 Firewall}

The A6 firewall is the structural boundary preventing non-dual algebra from
leaking downward into the dual-algebra region. It enforces continuity,
spectral stability, holonomy bounds, collapse locality, semantic lift, and
kernel preservation. These constraints ensure that diagnostic operators
(such as the Boundary Mapper) can safely read the manifold without inducing
geometric fracture.

\subsection{Firewall Purpose}

The firewall serves four primary functions:

\begin{enumerate}
    \item \textbf{Geometric Preservation:} Maintain curvature continuity and
          prevent geometric tearing.
    \item \textbf{Spectral Stability:} Ensure operator spectra vary
          continuously and remain bounded.
    \item \textbf{Holonomy Control:} Prevent drift-field accumulation from
          exceeding fracture thresholds.
    \item \textbf{Invariant Lift:} Preserve modal, semantic, and provenance
          invariants across operator application.
\end{enumerate}

These functions collectively define the safe region in which dual-algebra
operators remain reversible and diagnostically readable.

\subsection{Constructive Fixes}

The firewall is implemented through thirteen constructive fixes. Each fix
defines a boundary condition that must hold throughout the A6 manifold.

\paragraph{Fix 1: Curvature Continuity}
\[
K(p) \approx K(q) \quad \text{for } p \sim q.
\]

\paragraph{Fix 2: Cosine-Law Curvature Combination}
\[
K_{A\oplus B}
  = \sqrt{K_A^2 + K_B^2 + 2 K_A K_B \Phi(A,B)}.
\]

\paragraph{Fix 3: Drift-Field Interaction}
\[
[\nabla_X,\nabla_Y] = \nabla_{[X,Y]}.
\]

\paragraph{Fix 4: Modal Adjacency}
\[
\mu(p^+) \in \mathcal{N}(\mu(p^-)).
\]

\paragraph{Fix 5: Spectral Continuity}
\[
\mathrm{Spec}(L(p)) \text{ continuous in } p.
\]

\paragraph{Fix 6: Spectral Mismatch Bound}
\[
\delta_{\mathrm{spec}}(p) < \epsilon.
\]

\paragraph{Fix 7: Holonomy Drift Bound}
\[
H(\gamma) < \Theta_{\max}.
\]

\paragraph{Fix 8: Collapse-Surface Locality}
\[
p \in \Sigma \Rightarrow \exists \epsilon > 0 : B_\epsilon(p) \cap \Sigma \neq \emptyset.
\]

\paragraph{Fix 9: Stability Envelope Monotonicity}
\[
E(p) \le E(q) \Rightarrow \text{Stability}(p) \Rightarrow \text{Stability}(q).
\]

\paragraph{Fix 10: Modal Lift Preservation}
\[
\mu(p) \leadsto \mu(q).
\]

\paragraph{Fix 11: Provenance Lift Preservation}
\[
I_{\mathrm{prov}}(M) \Rightarrow I_{\mathrm{prov}}(\mathcal{A}(M,S)).
\]

\paragraph{Fix 12: Semantic Lift Preservation}
\[
I_{\mathrm{sem}}(M) \Rightarrow I_{\mathrm{sem}}(\mathcal{A}(M,S)).
\]

\paragraph{Fix 13: Operator Determinism and Kernel Preservation}
\[
\mathcal{O}(p) = \mathcal{O}(q) \Rightarrow p = q,
\]
\[
K_{\mathrm{A6}}(p) \Rightarrow K_{\mathrm{A6}}(\mathcal{O}(p)).
\]

\subsection{Firewall Summary}

The A6 firewall defines the structural constraints required for safe
diagnostic operation. It ensures that the Boundary Mapper can read the
manifold without inducing collapse, drift fracture, spectral tearing, or
semantic discontinuity. The firewall is not a diagnostic tool; it is the
mathematical substrate that makes diagnostics possible.
